\section{Conclusion}
\label{sec:conclusion}

This survey has provided a comprehensive review of tactile sensing for dexterous robot manipulation, spanning sensor technologies, hand design, data representation, learning methods, VLA models, and embodiment retargeting. Our analysis of 146 papers (2020--2026) reveals three converging trends that are driving the field at unprecedented pace: (1) cost reduction of vision-based sensors (BioTac \$10K $\to$ DIGIT \$350 $\to$ ReSkin \$6), (2) democratization of open-source hands (Shadow \$100K $\to$ LEAP \$2K), and (3) emergence of VLA foundation models (RT-2 $\to$ pi0 $\to$ Gemini Robotics) with nascent tactile integration (ForceVLA, Tactile-VLA).

We identified the top 10 open challenges, led by the tactile sim-to-real gap and the absence of general solutions for cross-embodiment tactile transfer. Current industrial deployments remain at the logistics level; dexterous assembly has not yet reached production floors. To bridge this gap, we proposed the mechanism--tactile--learning triangle, which combines physical intelligence (underactuation, VSA, active surfaces) with tactile perception and learned policies to reduce the burden on each individual axis.

As of 2026, tactile sensing is transitioning from optional to standard in humanoid robotics. Sanctuary AI has integrated 5~mN sensitivity into a commercial platform, NVIDIA generates 780K synthetic trajectories in 11 hours, and pi0.5/RECAP demonstrates continuous post-deployment improvement. The question is no longer \emph{whether} robots need touch, but \emph{how quickly} the field can deliver manufacturing-grade dexterous manipulation. We hope this survey accelerates that transition.
